\section{Action 1.7 : Câblage des Références Inter-Projets (ProjectReference)}

\begin{tcolorbox}[colback=indigo!4!white,colframe=indigo!40!white,arc=2mm,boxrule=0.8pt,top=2mm,bottom=2mm,left=3mm,right=3mm]
    \small\textbf{\color{indigo!90!black}Objectif d'apprentissage} : À l'issue de cette fiche, l'apprenant sera capable de distinguer le regroupement par solution de la liaison binaire par \texttt{ProjectReference}, de relier les projets exécutables (\texttt{Back} et \texttt{Front}) à la bibliothèque de contrats (\texttt{Shared}) et d'expliquer pourquoi une référence directe entre Front et Back est strictement interdite.
\end{tcolorbox}

\vspace{0.1cm}

Au terme de l'Action 1.5, les trois projets coexistent au sein du conteneur \texttt{FullStackCSharp.sln}. Cependant, sur le plan du compilateur C\#, ils sont encore totalement étrangers les uns aux autres : le serveur \texttt{Back} et le client \texttt{Front} sont incapables de consommer les futurs DTOs hébergés dans \texttt{Shared}.

L'\textbf{Action 1.7} concrétise la liaison binaire officielle en ajoutant une balise \textbf{\texttt{<ProjectReference>}} dans les fichiers \texttt{.csproj} respectifs :
\begin{itemize}
    \item \texttt{FullStackCSharp.Back.csproj} $\longrightarrow$ référence \texttt{FullStackCSharp.Shared.csproj}
    \item \texttt{FullStackCSharp.Front.csproj} $\longrightarrow$ référence \texttt{FullStackCSharp.Shared.csproj}
\end{itemize}

\vspace{0.3cm}

\begin{schemabox}{Topologie des Références de Compilation (L'Autorité des Contrats)}
\begin{center}
\begin{tikzpicture}[
    boxshared/.style={
        rectangle, 
        draw=cyanproj!90!black, 
        fill=cyanproj!7!white, 
        rounded corners=3mm, 
        minimum width=8.6cm, 
        minimum height=1.7cm, 
        inner ysep=2.5mm,
        align=center, 
        line width=1.1pt
    },
    boxapp/.style={
        rectangle, 
        draw=#1!80!black, 
        fill=#1!6!white, 
        rounded corners=2.5mm, 
        minimum width=4.8cm, 
        minimum height=2.1cm, 
        inner ysep=3mm,
        align=center, 
        line width=0.9pt
    },
    refarrow/.style={
        -{Stealth[length=3mm]}, 
        very thick, 
        color=cyanproj!90!black
    },
    forbidbox/.style={
        rectangle, 
        draw=red!70!black, 
        fill=red!5!white, 
        rounded corners=2mm, 
        minimum width=3.8cm, 
        minimum height=0.9cm, 
        inner ysep=2mm,
        align=center, 
        line width=0.8pt,
        densely dashed
    },
    forbidarrow/.style={
        thick, 
        densely dashed, 
        color=red!70!black
    }
]

% Top: Shared Contract Library (y = 3.6)
\node[boxshared] (shared) at (0, 3.6) {
    \textbf{\large FullStackCSharp.Shared}\\[2pt]
    {\color{cyanproj!80!black}\rule{7.2cm}{0.4pt}}\\[2pt]
    \footnotesize \textbf{Autorité Centrale des Modèles d'Échange (DTOs)}\\
    \tiny Produit binaire : \texttt{FullStackCSharp.Shared.dll}
};

% Bottom: Front and Back (y = 0.4)
\node[boxapp=indigo] (front) at (-5.2, 0.4) {
    \textbf{\normalsize FullStackCSharp.Front}\\[2pt]
    {\color{indigo!70!black}\rule{3.8cm}{0.4pt}}\\[2pt]
    \scriptsize Client Mobile \& Desktop (.NET MAUI)\\[2pt]
    \tiny Dépendance : \texttt{<ProjectReference>} vers Shared\\
    \tiny Accède aux types de requêtes / réponses
};

\node[boxapp=secondary] (back) at (5.2, 0.4) {
    \textbf{\normalsize FullStackCSharp.Back}\\[2pt]
    {\color{secondary!70!black}\rule{3.8cm}{0.4pt}}\\[2pt]
    \scriptsize Serveur ASP.NET Core (Kestrel)\\[2pt]
    \tiny Dépendance : \texttt{<ProjectReference>} vers Shared\\
    \tiny Mappe les DTOs vers les endpoints
};

% ProjectReference curves entering Shared cleanly from sides
\draw[refarrow] (front.north) to[out=90, in=180] 
    node[midway, left=4pt, font=\scriptsize\bfseries\color{cyanproj!90!black}, align=right] {ProjectReference\\\texttt{Shared.csproj}}
    (shared.west);

\draw[refarrow] (back.north) to[out=90, in=0] 
    node[midway, right=4pt, font=\scriptsize\bfseries\color{cyanproj!90!black}, align=left] {ProjectReference\\\texttt{Shared.csproj}}
    (shared.east);

% Middle: Forbidden Direct Coupling
\node[forbidbox] (forbid) at (0, 0.4) {
    \scriptsize\bfseries\color{red!80!black} $\times$~Liaison Directe Interdite\\[1pt]
    \tiny Front et Back s'ignorent mutuellement
};
\draw[forbidarrow] (-2.7, 0.4) -- (forbid.west);
\draw[forbidarrow] (forbid.east) -- (2.7, 0.4);

\end{tikzpicture}
\end{center}
\end{schemabox}

\vspace{0.2cm}

\begin{bestpracticebox}{L'Interdiction Absolue du Couplage Direct Front $\leftrightarrow$ Back}
Dans une architecture Fullstack moderne et sécurisée :
\begin{itemize}
    \item \textbf{Le client MAUI ne doit JAMAIS référencer le projet Back} : Référencer le serveur exposerait l'application cliente aux entités internes de la base de données SQLite centrale et aux secrets d'infrastructure (chaînes de connexion, clés privées JWT).
    \item \textbf{Le serveur Kestrel ne doit JAMAIS référencer le projet Front} : Le serveur est un processus hébergé sur une machine distante ; il n'a aucune raison d'embarquer les contrôles d'interface graphique XAML.
    \item \textbf{Le seul pont autorisé est le réseau HTTP} : Front et Back communiquent exclusivement par le protocole HTTP transportant des données JSON conformes aux contrats définis dans \texttt{Shared}.
\end{itemize}
\end{bestpracticebox}

\newpage

\subsection*{Manipulation Reproductible : Câblage des Références}

\begin{bestpracticebox}{Ajout des Références via la CLI}
\textbf{Répertoire de travail} : Racine du dépôt (\texttt{DEV/FullStackCSharp}).\\[2pt]
\textbf{Commandes CLI} :
\begin{verbatim}
dotnet add FullStackCSharp.Back reference FullStackCSharp.Shared
dotnet add FullStackCSharp.Front reference FullStackCSharp.Shared
\end{verbatim}
\textbf{Résultat attendu} : Chaque fichier \texttt{.csproj} enrichit son bloc \texttt{<ItemGroup>} d'une balise :
\begin{verbatim}
<ProjectReference Include="..\FullStackCSharp.Shared\FullStackCSharp.Shared.csproj" />
\end{verbatim}
\textbf{Vérification concrète} : Lancer la compilation globale :\\
\texttt{dotnet build FullStackCSharp.sln} $\longrightarrow$ MSBuild compile automatiquement \texttt{Shared} en premier, puis compile \texttt{Back} et \texttt{Front} en utilisant l'assembly binaire généré.
\end{bestpracticebox}

\vspace{0.2cm}

\subsection*{Erreurs Fréquentes du Débutant}

\begin{itemize}
    \item \textbf{Confondre \texttt{dotnet sln add} et \texttt{dotnet add reference}} : Le premier enregistre le projet dans le classeur de la solution ; le second autorise un projet à utiliser le code et les types d'un autre projet.
    \item \textbf{Créer une référence circulaire (A dépend de B et B dépend de A)} : C'est impossible dans .NET et rejeté immédiatement par MSBuild avec l'erreur \texttt{Circular dependency detected}.
\end{itemize}

\vspace{0.2cm}

\begin{tcolorbox}[colback=emerald!4!white,colframe=emerald!40!white,arc=2mm,boxrule=0.8pt,top=2.5mm,bottom=2.5mm,left=3mm,right=3mm]
    \small\textbf{\color{emerald!80!black}Exercice de réinvestissement} :
    Ouvrez le fichier \texttt{FullStackCSharp.Back/FullStackCSharp.Back.csproj}. Identifiez la balise \texttt{<ProjectReference>}. Quelle est la syntaxe du chemin relatif vers \texttt{Shared} ?\\
    \textbf{Critère d'auto-évaluation} : Vous devez identifier le chemin relatif : \texttt{..\textbackslash FullStackCSharp.Shared\textbackslash FullStackCSharp.Shared.csproj}.
\end{tcolorbox}

\vspace{0.3cm}

\subsection*{Questions Fréquentes du Débutant}

\faqitem{Pourquoi le chemin utilise-t-il \texttt{..} (deux points) ?}{
    Le fichier \texttt{FullStackCSharp.Back.csproj} est situé dans le sous-dossier \texttt{FullStackCSharp.Back/}. Le préfixe \texttt{..} indique à MSBuild de « remonter d'un cran » vers le dossier parent (la racine de la solution) pour ensuite redescendre dans le sous-dossier \texttt{FullStackCSharp.Shared/}.
}

\faqitem{Que se passe-t-il si je modifie un fichier dans \texttt{Shared} et que je lance \texttt{dotnet build FullStackCSharp.Back} ?}{
    Grâce à la balise \texttt{<ProjectReference>}, MSBuild inspecte les dates de modification de tous les projets de la chaîne de dépendance. Il détecte automatiquement que \texttt{Shared} a été altéré, le recompile en priorité vers une nouvelle DLL fraîche, puis compile \texttt{Back}. Vous n'avez jamais besoin de recompiler manuellement les dépendances intermédiaires.
}
